%% ============================================================
%%  Section: Main Results
%%  Depositor Discipline in Modern Banking (2026-04-03 draft)
%% ============================================================

\section{Main Results}
\label{sec:results}

We now test whether the CECL Day-One disclosure triggered a disciplinary
response from uninsured depositors. Our design isolates three empirical
predictions implied by the theory: (i) banks whose adoption-quarter CECL
adjustment revealed larger latent credit losses should experience relative
declines in uninsured deposit balances after the disclosure; (ii) the shift
should be concentrated in the uninsured portion of the deposit portfolio and not
replicated symmetrically in insured balances; and (iii) the informational cost of
losing uninsured funding should translate into higher funding costs and lower
profitability, the hallmark of a genuine disciplinary mechanism rather than a
mechanical accounting effect. The results confirm all three predictions. We
proceed by first establishing the quantity response in time deposits, then
confirming it with a triple-difference specification that absorbs bank-wide
funding trends, and finally documenting the transmission to bank financial
performance.

\subsection{Depositor Quantity Response: Time Deposits}
\label{sec:results:time_deposits}

Table~\ref{tab:small_bank_did_time_deposits} presents difference-in-differences
estimates for uninsured and insured time deposits, separately, over the 97,064
bank-quarter observations in the community-bank panel. All specifications absorb
bank and calendar-quarter fixed effects and include lagged log assets,
equity-to-assets, and interest-expense-to-assets as controls.

\paragraph{Uninsured time deposits.}
The interaction of the post-adoption indicator with the binary high-CECL
treatment (column~2) yields a coefficient of $-0.394$ (standard error $0.182$),
significant at the 5 percent level. Relative to otherwise similar low-CECL
banks, high-CECL community banks saw their uninsured time deposits shrink by
roughly 39 basis points of total assets in the post-adoption period. This
magnitude is economically meaningful. The median community bank in our sample
holds approximately \$500 million in total assets, so a 39-basis-point
contraction in uninsured time deposits implies roughly \$2 million in lost
uninsured balances per bank. Scaled to the top decile of the CECL distribution,
where the average adjustment exceeds 3.2 percent of pre-adoption equity, the
implied balance contraction represents a significant fraction of a typical
community bank's uninsured funding base.

The continuous-treatment estimate (column~1) is consistent: a one-percentage-point
increase in the CECL equity shock reduces uninsured time deposits by $0.072$
percentage points of assets ($p < 0.05$). The two estimates are directly
comparable—at the top-decile threshold the continuous specification predicts a
contraction of approximately 23 basis points, reasonably close to the binary
estimate of 39 basis points given that the high-CECL indicator averages across
banks just above the threshold.

\paragraph{Insured time deposits.}
The pattern for insured time deposits runs in the opposite direction. Column~4
reports a high-CECL coefficient of $+0.712$ ($p < 0.05$), indicating that
high-exposure banks attracted substantially more insured time deposits after
adoption. The contrast with the uninsured estimate is economically striking: the
same banks that shed 39 basis points of uninsured balances simultaneously added
71 basis points of insured balances, a net increase in total time deposits. The
continuous insured estimate in column~3 is positive ($+0.070$) but does not
reach conventional significance thresholds.

This asymmetric response is not incidental—it is the disciplinary mechanism in
operation. Uninsured depositors, exposed to loss in the event of bank failure,
withdrew funds upon observing that their bank carried higher-than-expected latent
credit losses. High-CECL banks, facing the loss of this funding, had to
substitute toward insured deposits, bidding up rates on certificates of deposit
to attract replacement balances. The result is a compositional shift in the
deposit portfolio toward the more expensive, insurance-protected tier—precisely
the dynamic that theory predicts when informed creditors withdraw and banks must
replace them with uninformed, insured ones. We document the cost consequences
of this substitution in Section~\ref{sec:results:financial}.

\paragraph{Event-study dynamics and pre-trends.}
Figure~\ref{fig:es_time_deps_cecl} plots event-study coefficients for the
window of ten quarters on either side of CECL adoption ($k = -1$ omitted).
Two features are diagnostic. First, for both the binary (left panels) and
continuous (right panels) treatment, the pre-adoption coefficients for uninsured
time deposits cluster tightly around zero through all ten lead quarters and are
jointly statistically indistinguishable from the reference period. High-CECL and
low-CECL banks were on essentially identical deposit trajectories before the
disclosure—a clean validation of the parallel-trends assumption. Second, the
post-adoption break is sharp: the uninsured coefficient turns negative within the
first adoption quarter and remains persistently below zero, while the insured
coefficient moves in the opposite direction with comparable persistence. The
symmetry in timing and the absence of anticipation effects reinforce the
interpretation of the disclosure as the operative trigger.

\subsection{Triple-Difference: Isolating Within-Bank Composition}
\label{sec:results:triple}

A natural concern is that high-CECL banks may have faced general funding
pressures after 2023 that shifted their entire time-deposit portfolio
independently of depositor discipline. To address this, Table~\ref{tab:small_bank_triple_diff_time_deposits}
stacks the uninsured and insured time-deposit series in a single regression and
estimates a triple difference: whether the \textit{relative} change in uninsured
versus insured deposits is differentially larger for high-CECL banks
post-adoption. The specification includes the full set of pairwise interactions,
bank and quarter fixed effects, and the same lagged controls.

The triple interaction---$\text{High CECL} \times \text{Post} \times
\text{Uninsured}$ in column~(2)---is $-1.252$ (standard error $0.340$),
significant at the 1 percent level. For the continuous specification (column~1),
the analogous triple interaction is $-0.177$ (standard error $0.058$), also
significant at the 1 percent level. These estimates are identified from
within-bank variation in the composition of insured versus uninsured time
deposits before and after adoption, controlling for any bank-wide post-adoption
funding trend through the $\text{High CECL} \times \text{Post}$ two-way
interaction. The triple-difference design therefore rules out the alternative
hypothesis that CECL simply depressed aggregate deposit demand at high-exposure
banks: if that were the explanation, insured and uninsured balances would move
together, and the triple interaction would be zero.

The magnitude is large. High-CECL banks experienced a 125-basis-point widening
of the gap between their insured and uninsured time-deposit balances (as a share
of assets) relative to the pre-adoption composition and relative to low-CECL
banks. This gap materializes through both legs simultaneously: uninsured deposits
fall and insured deposits rise, so the composition shifts rapidly toward the
protected tier. The $\text{Post} \times \text{Uninsured}$ coefficient in
column~(2) is $+0.898$ ($p < 0.01$), confirming that across the entire
community-bank sample uninsured time deposits also rose after 2023Q1—plausibly
reflecting the broader post-pandemic rotation of deposits into higher-yielding
time accounts as the Federal Reserve's tightening cycle raised CD rates. The
triple interaction with High CECL captures the \textit{differential} within-bank
reallocation away from uninsured deposits at the banks whose credit-risk
information was revised most adversely by CECL adoption, above and beyond this
common trend.

Figure~\ref{fig:triple_diff_time_deps_es} plots the corresponding event-study
coefficients for the triple interaction. Pre-adoption leads are uniformly flat in
both the binary (Panel A) and continuous (Panel B) specifications. The triple
interaction turns sharply negative beginning at adoption and widens for two to
three quarters before stabilizing, suggesting that the composition adjustment was
largely complete within the first year of disclosure.

\subsection{Financial Consequences: Funding Costs and Profitability}
\label{sec:results:financial}

Table~\ref{tab:small_bank_did_financial_outcomes} turns to the financial
consequences of the deposit reallocation. The dependent variables are quarterly
interest expense scaled by assets (columns~1--2), return on assets (columns~3--4),
and net interest margin (columns~5--6). All specifications include bank and
quarter fixed effects and lagged log assets and equity ratio.

\paragraph{Funding costs.}
Column~(2) shows that high-CECL banks paid $0.021$ percentage points more in
quarterly interest expense per dollar of assets after adoption ($p < 0.01$).
Annualized, this amounts to approximately 8--9 basis points of additional funding
cost relative to low-CECL peers. The continuous estimate in column~(1)
implies that a one-percentage-point increase in the CECL equity shock raises
quarterly interest expense by 0.37 basis points of assets, consistent with the
binary estimate scaled at the top-decile threshold.

The mechanism is straightforward: having lost lower-cost uninsured deposits,
high-CECL banks replaced them with insured time deposits at rates sufficient to
attract depositors who would not otherwise bear uninsured credit risk at these
institutions. The cost differential is the premium community banks paid for their
disciplined behavior being revealed by CECL.

\paragraph{Profitability.}
The profitability effects confirm that the funding cost increase was not absorbed
by higher loan yields or cost-cutting elsewhere. Return on assets fell by $0.038$
percentage points per quarter for high-CECL banks (column~4, $p < 0.01$), an
annualized drag of roughly 15 basis points. For community banks operating at
typical ROA levels of 100--120 basis points annually in 2023--2024, this
represents a non-trivial compression of about 12--15 percent of average earnings.
Net interest margin also contracted by $0.024$ percentage points per quarter
(column~6, $p < 0.05$), an annualized effect of approximately 9 basis points.
The joint decline in ROA and NIM confirms that the higher interest expense on
the liability side was not passed through to borrowers—consistent with a
competitive loan market in which community banks are price-takers on assets but
must bid up liabilities to compensate for the loss of uninformed uninsured
depositors.

The continuous treatment estimates scale proportionally: each percentage point of
CECL equity shock reduces quarterly ROA by $0.009$ percentage points ($p < 0.01$)
and quarterly NIM by $0.005$ percentage points ($p < 0.05$), reinforcing that
the effects are a smooth function of disclosure severity rather than a threshold
phenomenon concentrated at the very top of the CECL distribution.

\paragraph{Event-study dynamics.}
Figure~\ref{fig:es_financial_outcomes_cecl} plots event-study coefficients for
all three financial outcomes. Pre-adoption leads are flat for interest expense,
ROA, and NIM across both treatment specifications, again supporting parallel
trends. The post-adoption divergence for ROA and NIM is persistent through the
end of the sample horizon—consistent with a multi-year repricing dynamic in
which maturing insured CDs, raised to replace the outflowing uninsured balances,
continue to roll over at elevated rates in subsequent periods. The interest
expense effect also persists, indicating that the cost of disciplinary funding
substitution is not merely a transitional adjustment but reflects a durable
repricing of community bank liabilities at high-CECL institutions.

\subsection{Summary}
\label{sec:results:summary}

Taken together, the evidence in Tables~\ref{tab:small_bank_did_time_deposits}--\ref{tab:small_bank_did_financial_outcomes}
and Figures~\ref{fig:es_time_deps_cecl}--\ref{fig:es_financial_outcomes_cecl}
supports a coherent depositor discipline narrative triggered by the CECL
information shock. At the disclosure date, high-exposure community banks
revealed that their loan portfolios carried materially larger unrecognized
expected credit losses than outsiders had previously estimated. Uninsured
depositors responded by reducing their balances, exploiting the liquidity of
demand and time deposits to act on the new information. High-CECL banks
replaced the outflow by bidding up rates on insured time deposits—the only
funding tier not subject to run risk once risk is revealed—paying an enduring
premium in funding costs that compressed both net interest margins and return on
assets. The triple-difference specification confirms that this is a within-bank
composition shift in the structure of the deposit portfolio rather than a
uniform funding retreat, and parallel pre-trends in all event-study plots
support a causal interpretation. The results establish that uninsured depositor
discipline operates in the non-crisis, non-TBTF setting where governance theory
predicts it should be strongest—at small community banks, far from the implicit
backstops that have attenuated evidence of the mechanism elsewhere.
